\section{Introduction}\label{sec:introduction}

Factor models represent systematic return variation through exposures to
common shocks, and the resulting covariances determine how firm risks
aggregate in a portfolio
\citep{ross_arbitrage_1976,fama_multifactor_1996,daniel_evidence_1997}.
Yet the characteristics used to describe those exposures are usually encoded
as points or finite vectors, even when the evidence about a firm consists of
heterogeneous observations measured with uncertainty.
Such a point representation does not itself encode how the characteristic
evidence is distributed within a firm.
We therefore ask what restrictions become available when firm characteristics
are treated instead as probability distributions.

Under this representation, each firm is a cloud of characteristic
observations rather than a single location.
The distance between two clouds reflects how far their probability mass must
be moved to align them.
A coupling is a joint arrangement that specifies which draws from the two
distributions occur together.
Quadratic Wasserstein distance, denoted W$_2$, formalizes this comparison by
selecting the coupling that minimizes expected squared displacement between
the two distributions.
The covariance result applies the same geometry after characteristics have
been mapped into latent factor exposures.
Given two latent exposure laws, quadratic Wasserstein distance in factor-risk
coordinates yields sharp upper and lower covariance endpoints over admissible couplings.
Polarization expresses covariance as average exposure magnitudes minus half
the expected squared distance, so the least-cost coupling supplies the
covariance ceiling and reflection supplies the floor.
A maintained common transmission map then transfers a one-sided restriction
on the covariance ceiling to characteristic-side distance.
Any realized coupling then incurs a nonnegative transport-excess cost
relative to the maximum-covariance arrangement.

This result is conditional rather than a point prediction.
It requires a common randomized bi-Lipschitz transmission map, bounds on
firm-level slack expressed in risk coordinates, and a maintained bridge
between factor exposures and return covariance.
The article and return data do not identify the map, its constants, or the
latent coupling.
Accordingly, the contribution concerns theory and measurement of covariance,
not expected returns or no-arbitrage pricing.

\Cref{fig:continuous-bound-construction} summarizes the construction from
article-embedding proxies for characteristic distributions to a conditional
interval for the covariance ceiling.
The observed articles proxy each firm's characteristic distribution, while
the exposure laws and transmission restrictions remain latent.

\begin{figure}[H]
  \centering
  \ppfigure{1}{continuous_bound_construction}
  \caption{Construction of the one-sided systematic-covariance envelope.
    Panel A maps each firm's characteristic law through the common
    transmission map $T$ with kernel $K$ to its latent exposure law $P_i$.
    Panel B states the bi-Lipschitz and risk-coordinate slack restrictions on
    the shared randomized map.
    Panel C combines the exposure-law covariance ceiling with transfer bounds
    to form the conditional interval evaluated from observed articles.
  The operator $\Gamma$ fixes the common risk-coordinate geometry.}
  \label{fig:continuous-bound-construction}
\end{figure}

Empirically, we approximate each firm's latent characteristic distribution by
the empirical distribution of language-model representations of financial
news articles, disclosures, and analyst reports.
These firm-level article clouds come from a Nasdaq sample spanning
\ppvalue{sample-year-start}--\ppvalue{sample-year-end}.
We then relate pairwise article-embedding W$_2$ distance to return
co-movement in a dyadic regression with symmetric firm effects.
The canonical association is positive: one standard deviation more W$_2$
corresponds to roughly eight fewer correlation points.
A timing-separated exercise holds the article-embedding geometry fixed and
recomputes return outcomes in a later window, for which the estimated
association retains the same sign.
This comparison is corroborative rather than point-in-time predictive evidence.
The regression coefficient is associational, not causal, and is not a
numerical estimate of transport excess.

A separate sensitivity grid evaluates how informative the conditional
envelope remains as its maintained restrictions are relaxed.
In that grid, $L$ governs multiplicative distortion in the
characteristic-to-exposure map and $\tau$ governs additive slack in risk coordinates.
A covered dyad has nonnegative transport excess throughout its
scenario-implied interval, a violated dyad has negative excess throughout,
and an unresolved dyad has an interval spanning zero.
The zero-slack isometric cell $(L,\tau)=(1,0)$ classifies
\ppnum[3]{confound-w2-l1-tau0-coverage} of observed dyads as covered and
\ppnum[3]{confound-w2-l1-tau0-violation} as violations.
Even $(L,\tau)=(1,0.05)$ and $(1,0.10)$ leave
\ppnum[3]{confound-w2-l1-tau0p05-unresolved} and
\ppnum[3]{confound-w2-l1-tau0p1-unresolved} unresolved, respectively.
The envelope is therefore informative only when the
characteristic-to-exposure map is tightly restricted.

The paper makes three connected contributions.
First, it derives sharp pairwise covariance endpoints for latent exposure
laws and transfers a conditional restriction on the covariance ceiling to
distribution-valued characteristics.
Second, it provides a measurement design that keeps the distinction between
latent characteristic laws, article-based proxies, and return covariance explicit.
Third, it reports both the reduced-form distance association and the scenario
evidence needed to assess when the theoretical envelope is informative.

\Cref{sec:literature} positions the covariance contribution relative to
characteristic-based loading models and textual proximity measures.
\Cref{sec:methodology} constructs the one-sided covariance envelope and
states its maintained conditions.
\Cref{sec:empirical-design,sec:results} develops the proxy and estimation
design before reporting the association, timing, and sensitivity evidence.
